% GENERATED FILE. Edit canonical lessons, then run npm run generate:books.
% canonical-source-hash: fnv1a64:c4d2310777bd3178
% canonical-lessons: LA-C143-durus, LA-C143-mollis, LA-C143-tepidus, LA-C143-antiquus, LA-C143-angustus

\chapter{Hard, Soft, Lukewarm, Ancient, Narrow}
\label{ch:la-hard-soft-lukewarm-ancient-narrow}
\hlchaptermodality{143}

\begin{chapteropening}
\textbf{By the end of this chapter:} \emph{I can say hard, soft, lukewarm, ancient and narrow.}

Complete the last lesson of chapter 143: I can say hard, soft, lukewarm, ancient and narrow.
\end{chapteropening}

\section[dūrus, dūra, dūrum]{dūrus, dūra, dūrum — hard}

\label{lesson:LA-C143-durus}



\begin{quote}
Before the new one: say the Latin for a flute, then the Latin for a lyre.
\end{quote}

\subsection*{You'll want to know: dūrus, dūra, dūrum}
\textbf{dūrus, dūra, dūrum} — \textquotedblleft{}hard\textquotedblright{}.

\begin{cousinweb}
\textbf{dūrus} has no sure origin; Lewis \& Short call it doubtful. A link with Proto-Indo-European \emph{*deru-}, \textquotedblleft{}firm, solid\textquotedblright{}, the root of words for tree, is often proposed (EtymDB 2.1).
\end{cousinweb}

\begin{grammarlens}[title={What you've built}]
One more word for this chapter.
\end{grammarlens}

\subsection*{Your turn}
\begin{itemize}
\raggedright
  \item \emph{Say it:} \emph{dūrus}
  \item \emph{Say it:} \emph{dūrus}, once more
  \item \emph{Say it:} say \emph{dūrus}
  \item \emph{Recall:} say the Latin for a flute, then the Latin for a lyre, then say \emph{dūrus} again
\end{itemize}

\begin{tcolorbox}[breakable,colback=teal!4,colframe=teal!35!black,title={Before you move on}]
What does dūrus, dūra, dūrum mean? (\textquotedblleft{}Hard\textquotedblright{}.) Say it once more.
\end{tcolorbox}
\section[mollis, molle]{mollis, molle — soft}

\label{lesson:LA-C143-mollis}



\begin{quote}
Before the new one: say the Latin for a lyre, then the Latin for hard.
\end{quote}

\subsection*{You'll want to know: mollis, molle}
\textbf{mollis, molle} — \textquotedblleft{}soft\textquotedblright{}.

\begin{cousinweb}
\textbf{mollis} is inherited, usually traced to Proto-Indo-European \emph{*moldus}, \textquotedblleft{}soft, weak\textquotedblright{}, from \emph{*mel-}, \textquotedblleft{}soft\textquotedblright{} (EtymDB 2.1); the ld became ll. Greek \emph{amalos}, \textquotedblleft{}weak\textquotedblright{}, is compared (Lewis \& Short).
\end{cousinweb}

\begin{grammarlens}[title={What you've built}]
One more word for this chapter.
\end{grammarlens}

\subsection*{Your turn}
\begin{itemize}
\raggedright
  \item \emph{Say it:} \emph{mollis}
  \item \emph{Say it:} \emph{mollis}, once more
  \item \emph{Say it:} say \emph{mollis}
  \item \emph{Recall:} say the Latin for a lyre, then the Latin for hard, then say \emph{mollis} again
\end{itemize}

\begin{tcolorbox}[breakable,colback=teal!4,colframe=teal!35!black,title={Before you move on}]
What does mollis, molle mean? (\textquotedblleft{}Soft\textquotedblright{}.) Say it once more.
\end{tcolorbox}
\section[tepidus, tepida, tepidum]{tepidus, tepida, tepidum — lukewarm}

\label{lesson:LA-C143-tepidus}



\begin{quote}
Before the new one: say the Latin for hard, then the Latin for soft.
\end{quote}

\subsection*{You'll want to know: tepidus, tepida, tepidum}
\textbf{tepidus, tepida, tepidum} — \textquotedblleft{}lukewarm\textquotedblright{}.

\begin{cousinweb}
\textbf{tepidus} is built on \emph{tepeō}, \textquotedblleft{}to be warm\textquotedblright{} (Lewis \& Short), with \emph{-idus}, the ending of state adjectives. \emph{tepeō} is usually traced to Proto-Indo-European \emph{*tep-}, \textquotedblleft{}warm\textquotedblright{}.
\end{cousinweb}

\begin{grammarlens}[title={What you've built}]
One more word for this chapter.
\end{grammarlens}

\subsection*{Your turn}
\begin{itemize}
\raggedright
  \item \emph{Say it:} \emph{tepidus}
  \item \emph{Say it:} \emph{tepidus}, once more
  \item \emph{Say it:} say \emph{tepidus}
  \item \emph{Recall:} say the Latin for hard, then the Latin for soft, then say \emph{tepidus} again
\end{itemize}

\begin{tcolorbox}[breakable,colback=teal!4,colframe=teal!35!black,title={Before you move on}]
What does tepidus, tepida, tepidum mean? (\textquotedblleft{}Lukewarm\textquotedblright{}.) Say it once more.
\end{tcolorbox}
\section[antīquus, antīqua, antīquum]{antīquus, antīqua, antīquum — ancient, old-fashioned}

\label{lesson:LA-C143-antiquus}



\begin{quote}
Before the new one: say the Latin for soft, then the Latin for lukewarm.
\end{quote}

\subsection*{You'll want to know: antīquus, antīqua, antīquum}
\textbf{antīquus, antīqua, antīquum} — \textquotedblleft{}ancient, old-fashioned\textquotedblright{}.

\begin{cousinweb}
\textbf{antīquus} is built on \emph{ante}, \textquotedblleft{}before\textquotedblright{} (Lewis \& Short), and is usually traced to Proto-Indo-European \emph{*h2enti-h3kwos}, \textquotedblleft{}having been before\textquotedblright{} (EtymDB 2.1): what stood before, so old.
\end{cousinweb}

\begin{grammarlens}[title={What you've built}]
One more word for this chapter.
\end{grammarlens}

\subsection*{Your turn}
\begin{itemize}
\raggedright
  \item \emph{Say it:} \emph{antīquus}
  \item \emph{Say it:} \emph{antīquus}, once more
  \item \emph{Say it:} say \emph{antīquus}
  \item \emph{Recall:} say the Latin for soft, then the Latin for lukewarm, then say \emph{antīquus} again
\end{itemize}

\begin{tcolorbox}[breakable,colback=teal!4,colframe=teal!35!black,title={Before you move on}]
What does antīquus, antīqua, antīquum mean? (\textquotedblleft{}Ancient, old-fashioned\textquotedblright{}.) Say it once more.
\end{tcolorbox}
\section[angustus, angusta, angustum]{angustus, angusta, angustum — narrow}

\label{lesson:LA-C143-angustus}



\begin{quote}
Before the new one: say the Latin for lukewarm, then the Latin for ancient.
\end{quote}

\subsection*{You'll want to know: angustus, angusta, angustum}
\textbf{angustus, angusta, angustum} — \textquotedblleft{}narrow\textquotedblright{}.

\begin{cousinweb}
\textbf{angustus} is built on \emph{angō}, \textquotedblleft{}to squeeze, to choke\textquotedblright{} (Lewis \& Short), usually traced to Proto-Indo-European \emph{*h2engh-}, \textquotedblleft{}narrow, tight\textquotedblright{} (EtymDB 2.1). English \emph{anguish} and \emph{anxious} belong to the same family.
\end{cousinweb}

\begin{grammarlens}[title={What you've built}]
One more word for this chapter.
\end{grammarlens}

\subsection*{Your turn}
\begin{itemize}
\raggedright
  \item \emph{Say it:} \emph{angustus}
  \item \emph{Say it:} \emph{angustus}, once more
  \item \emph{Say it:} say \emph{angustus}
  \item \emph{Recall:} say the Latin for lukewarm, then the Latin for ancient, then say \emph{angustus} again
\end{itemize}

\begin{tcolorbox}[breakable,colback=teal!4,colframe=teal!35!black,title={Before you move on}]
What does angustus, angusta, angustum mean? (\textquotedblleft{}Narrow\textquotedblright{}.) Say it once more.
\end{tcolorbox}
